% Solutions to Chapter 6's exercises, from lectures/L06/appendix/d_solutions.md. None of the
% exercises is a program, so every one is answered.
%
% One correction to the source: 6.6 d) said that the bit lost when a register is counted up past
% 255 lands in the carry flag. That is true of an addition, but inc, which is how a register is
% counted up one step, leaves the carry flag alone (chapter 8). The source is corrected to match.

\answerchapter{c6:ch}

\answer{c6:ex:1}{Kombinatoriskt eller sekventiellt?}
\pt{a}\textbf{Kombinatoriskt.} Larmet beror bara på vad givarna visar just nu.

\pt{b}\textbf{Sekventiellt.} Med båda knapparna släppta kan reläet vara draget eller inte,
beroende på vilken knapp som trycktes sist. Kretsen minns.

\pt{c}\textbf{Kombinatoriskt.} Samma siffra in ger alltid samma segment tända.

\pt{d}\textbf{Sekventiellt.} Vilken färg som tänds härnäst beror på vilken som lyser nu.
Trafikljuset måste minnas sitt tillstånd.

\pt{e}\textbf{Sekventiellt.} Siffran 1 öppnar bara om 3 och 7 kom före. Låset måste minnas hur
långt i koden det har kommit.

\pt{f}\textbf{Kombinatoriskt.} Bromsbeslutet beror bara på de fyra ingångarna i det ögonblicket.

\medskip\noindent Frågan att ställa är alltid: \textbf{kan samma ingångar ge olika utgång?} Om ja,
finns ett minne.

\answer{c6:ex:2}{SR-låset}
\pt{a}\leavevmode

\begin{narrowtable}{@{}ccccl@{}}
  \thead{Steg} & \thead{S} & \thead{R} & \thead{Q} & \thead{Varför} \\
  \midrule
  start & 0 & 0 & 0 & givet \\
  1 & 1 & 0 & 1 & ettställ \\
  2 & 0 & 0 & 1 & minns \\
  3 & 0 & 0 & 1 & minns \\
  4 & 0 & 1 & 0 & nollställ \\
  5 & 0 & 0 & 0 & minns \\
  6 & 1 & 0 & 1 & ettställ \\
  7 & 0 & 0 & 1 & minns \\
\end{narrowtable}

\pt{b}Nej: steg 2, 3 och 7 har \code{Q = 1} och steg 5 har \code{Q = 0}, trots samma ingångar.
Utgången beror på vad som hände tidigare, och det är definitionen av ett sekvensnät. Ingen
sanningstabell med bara \code{S} och \code{R} som ingångar kan beskriva låset; tabellen måste också
ha med det tidigare \code{Q}.

\pt{c}Med \code{S = R = 1} har båda NOR-grindarna en etta på en ingång, så båda utgångarna blir 0.
\code{Q} och \code{Q'} är då inte längre varandras invers. Värre är vad som händer när båda släpps
samtidigt: båda grindarna ser då två nollor och vill bli 1, och den grind som hinner först vinner.
Resultatet går inte att förutsäga. En kombination som ger ett oförutsägbart minne är värdelös,
därför är den förbjuden.

\answer{c6:ex:3}{Lås och vippa}

\bookfigure{l06_sol_timing}{Klockan, \code{D} och utgångarna för lås och
vippa.}{c6:fig:sol_timing}

\pt{a}Låset följer \code{D} medan klockan är hög (de skuggade fönstren) och håller medan den är
låg: det blir 1 vid den första flanken, eftersom \code{D} redan är 1 då; 0 vid den andra, eftersom
\code{D} då är 0; 1 och sedan 0 igen under den korta pulsen, som kommer medan klockan fortfarande
är hög; oförändrat 0 genom det tredje fönstret; 1 vid den fjärde flanken, och kvar på 1.

\pt{b}Vippan läser \code{D} vid varje stigande flank: 1, 0, 0, 1, 1. Den enda skillnaden mot låset
är pulsen efter den andra flanken.

\pt{c}Pulsen kommer medan klockan är \textbf{hög}, strax efter den andra flanken. Då är låset
genomsläppligt, och pulsen går rakt igenom till låsets \code{Q}. Vippan ser den inte: den läste
\code{D} vid den andra flanken, innan pulsen kom, och läser nästa gång vid den tredje, när pulsen
redan är över. Hade pulsen kommit medan klockan var låg hade inte heller låset släppt igenom den.
Det är skillnaden mellan nivåkänsligt och flanktriggat, i ett enda exempel.

\ptnote{Vanligt fel}Att låta låset ändra sig bara vid flankerna, som vippan. Låset ändrar sig när
som helst medan klockan är hög, inte bara i början av det fönstret.

\answer{c6:ex:4}{Kontroll: bygg ett SR-lås}
\pt{a) och c}Tabellen från övning~\ref{c6:ex:2} ska stämma steg för steg. \code{Q'} ska vara
inversen av \code{Q} i varje steg, eftersom \code{S = R = 1} aldrig förekommer i följden. Om ett
steg skiljer sig är den vanligaste orsaken en felkopplad återkoppling: \code{Q} ska gå till den
\textbf{andra} grindens ingång, inte tillbaka till sin egen.

\pt{d}Innan någon ingång ändrats finns inget \enquote{rätt} värde. Kretsen har två stabila lägen,
\code{Q = 0} och \code{Q = 1}, och vilket den hamnar i beror på i vilken ordning simulatorn räknar
ut grindarna. En riktig krets hamnar på samma sätt i ett okänt läge när strömmen slås på, beroende
på små skillnader mellan transistorerna. Därför behöver varje verkligt system en \textbf{reset}: ett
sätt att ställa alla minnen i ett känt läge vid start. En mikrokontroller har en resetingång av
just det skälet.

\pt{e}Med \code{S = R = 1} blir både \code{Q} och \code{Q'} 0. När båda släpps samtidigt ser båda
grindarna två nollor, och i en simulator som räknar båda grindarna i samma steg kan de börja
\textbf{svänga}: båda blir 1, sedan båda 0, och så vidare. Andra gånger lägger sig kretsen i ett av
lägena. Båda utfallen bekräftar \secref{c6:a3}: kombinationen ger inget förutsägbart minne.

\answer{c6:ex:5}{Klockan}
\pt{a}\leavevmode
\[
  T = \frac{1}{16\,000\,000\ \text{Hz}} = 62{,}5\ \text{ns}
\]

\pt{b}16\,000\,000 cykler per sekund, alltså \textbf{16\,000 cykler per millisekund}.

\pt{c}1000 · 62,5~ns = 62\,500~ns = \textbf{62,5~µs}.

\pt{d}Halva frekvensen ger dubbla cykeltiden, 125~ns, och \textbf{125~µs}.

\answer{c6:ex:6}{Räknare}
\pt{a}En 3-bitars räknare börjar om efter 8 steg. 11 = 8 + 3, så den står på 3,
\textbf{\code{011}}.

\pt{b}999 kräver \textbf{10 bitar}: $2^9 = 512$ räcker inte, $2^{10} = 1024$ gör det. Tio vippor.

\pt{c}Varje bit har halva frekvensen av biten före: \code{Q0} 8~Hz, \code{Q1} 4~Hz, \code{Q2} 2~Hz
och \textbf{\code{Q3} 1~Hz}. \code{Q3} är en fyrkantvåg som är hög en halv sekund och låg en halv
sekund.

\pt{d}Det står på \textbf{0}. 255 är \code{1111 1111}, och ett till ger \code{1 0000 0000}, men den
nionde biten finns inte i ett 8-bitars register. Den tappas bort, och kvar blir \code{0000 0000}.
Det kallas rundgång. När mikrodatorn adderar hamnar den bortfallna biten i carry-flaggan, men inte
när den räknar upp med \code{inc}: \code{inc} lämnar carry-flaggan orörd (\secref{c8:a3}).

\answer{c6:ex:7}{Skiftregister}
\pt{a}\leavevmode

\begin{narrowtable}{@{}cccccc@{}}
  \thead{Efter flank} & \thead{in} & \thead{Q0} & \thead{Q1} & \thead{Q2} & \thead{Q3} \\
  \midrule
  start & & 0 & 0 & 0 & 0 \\
  1 & 1 & 1 & 0 & 0 & 0 \\
  2 & 1 & 1 & 1 & 0 & 0 \\
  3 & 0 & 0 & 1 & 1 & 0 \\
  4 & 1 & 1 & 0 & 1 & 1 \\
\end{narrowtable}

\pt{b}Före: \code{Q3 Q2 Q1 Q0 = 0011}. Vid flanken tar varje vippa värdet från vippan före, och
\code{Q0} får \code{in = 0}: \code{Q0 = 0}, \code{Q1 = 1}, \code{Q2 = 1}, \code{Q3 = 0}. Registret
innehåller \textbf{\code{0110}, talet 6}. Ett steg mot de mer signifikanta bitarna är
\textbf{multiplikation med 2}, precis som att lägga till en nolla sist i ett decimalt tal
multiplicerar med 10. Mikrodatorns instruktion \code{lsl} gör just det (kapitel~\ref{c8:ch}).

\answer{c6:ex:8}{Blockschemat}
\pt{a}Programräknaren, PC.
\pt{b}ALU:n.
\pt{c}Programminnet (flash).
\pt{d}Statusregistret, SREG, med noll-flaggan.
\pt{e}Styrenheten.
\pt{f}En utport, till exempel PORTB.
\pt{g}Dataminnet (SRAM).
\pt{h}Klockan.

\answer{c6:ex:9}{Instruktionscykeln}
\pt{a}\leavevmode

\begin{narrowtable}{@{}cclcccc@{}}
  \thead{Steg} & \thead{PC före} & \thead{Instruktion} & \thead{Händer} & \thead{r16} &
    \thead{r17} & \thead{PC efter} \\
  \midrule
  1 & 0 & \code{ldi r16, 10} & r16 = 10 & 10 & ? & 1 \\
  2 & 1 & \code{ldi r17, 20} & r17 = 20 & 10 & 20 & 2 \\
  3 & 2 & \code{add r16, r17} & r16 = 10 + 20 & 30 & 20 & 3 \\
  4 & 3 & \code{rjmp 3} & PC = 3 & 30 & 20 & 3 \\
  5 & 3 & \code{rjmp 3} & PC = 3 & 30 & 20 & 3 \\
  6 & 3 & \code{rjmp 3} & PC = 3 & 30 & 20 & 3 \\
\end{narrowtable}

\noindent\code{r17} har inget känt värde före steg 2; programmet har inte skrivit något där än.

\pt{b}Programmet hoppar till sig självt, om och om igen, så länge strömmen är på. En mikrokontroller
har inget operativsystem att återvända till när programmet är slut. Utan hoppet skulle PC fortsätta
till adress 4, 5, 6 och vidare, och CPU:n skulle utföra det som råkar ligga i flashminnet där, som
inte är ett program. En oändlig loop är hur ett program för en mikrokontroller slutar.

\answer{c6:ex:10}{Tre minnen}
\pt{a}Programmet i \textbf{flashminnet}. Knapptryckningsräknaren i \textbf{SRAM}.
Kalibreringsvärdet i \textbf{EEPROM}, som bevaras utan ström och kan skrivas en byte i taget.

\pt{b}Flash och EEPROM lagrar laddning på ett sätt som finns kvar utan matning. SRAM är i princip
lås (\secref{c6:a4}), och ett lås minns bara så länge dess grindar har ström: när matningen
försvinner försvinner återkopplingen, och därmed minnet.

\pt{c}Programminnet och dataminnet är \textbf{separata minnen med var sin buss}. CPU:n kan därför
hämta nästa instruktion samtidigt som den läser eller skriver data, och ett tal i programminnet nås
med andra instruktioner än ett tal i dataminnet.

\answer{c6:ex:11}{Bussar och minnesstorlek}
\pt{a}$2048 = 2^{11}$, så \textbf{11 adressbitar}.

\pt{b}32~kB = 32\,768 byte, och med 2 byte per instruktion blir det \textbf{16\,384
instruktioner}. $16\,384 = 2^{14}$, så programräknaren behöver \textbf{14 bitar}. (Några få
AVR-instruktioner är dubbelt så långa, 32 bitar, och tar två platser; det ändrar inte hur många
platser som finns.)

\pt{c}$2^8 =$ \textbf{256} olika värden, 0--255.

\answer{c6:ex:12}{Självhållning och SR-lås}
\pt{a}Startknappen \code{S1} ettställer, som \textbf{S}. Stoppknappen \code{S2} nollställer, som
\textbf{R}.

\pt{b}I reläkretsen sitter \code{S2} i serie med allt annat i spolens strömväg. När den trycks
bryts strömmen, vad \code{S1} än gör: \textbf{stopp vinner}. Kretsen kallas därför
\textbf{frånslagsdominerande}. I NOR-låset ger \code{S = R = 1} i stället två nollor på utgångarna
och ett oförutsägbart läge efteråt. Reläkretsen har alltså ingen förbjuden kombination: den har en
regel för vem som vinner, och regeln är vald för säkerhetens skull.

\pt{c}Låt \code{S1} gå förbi \code{S2}: koppla \code{S1} direkt från +24~V till spolen,
parallellt med seriekopplingen av \code{S2} och \code{K1}:s kontakt. Då drar \code{S1} alltid
reläet, även när \code{S2} är nedtryckt, och när \code{S1} släpps håller självhållningen genom
\code{K1} bara om \code{S2} inte är nedtryckt.

\begin{textdrawing}
            +24 V
       ┌──────┴──────┐
       │             │
   S2 (NC)       S1 (NO)
       │             │
   K1 (NO)           │
       └──────┬──────┘
         spole K1
              │
             0 V
\end{textdrawing}

\noindent Den varianten är ovanlig för maskiner, eftersom \textbf{stopp ska alltid vinna}: den som
trycker på stopp ska kunna lita på att maskinen stannar, även om någon annan samtidigt håller in
start eller en startknapp har fastnat. Nödstopp kopplas därför alltid så att de bryter, och de
bryter allt.
